\documentclass[10pt]{article}

\usepackage[utf8]{inputenc}

\usepackage[T1]{fontenc}

\usepackage{graphicx}

\usepackage[export]{adjustbox}

\graphicspath{ {./images/} }

\usepackage{amsmath}

\usepackage{amsfonts}

\usepackage{amssymb}

\usepackage[version=4]{mhchem}

\usepackage{stmaryrd}

\usepackage{mathrsfs}



\begin{document}

конечных элементов для уравнений с частными производными,



\includegraphics[max width=\textwidth, center]{2023_07_09_1c0a5893169b96c905ddg-1(1)}

с к и й А. А., Н и к о лі а в Е. С., Методы решения сеточных уравнений, M., 1978; [12] Д ж о p д ж А., Л Ю Д ж., Численное решение больших разреженных систем уравнений, пер. с англ., М., 1984; [13] Б р е б 6 и а K., у о к е p С.,

Применение метода граничных әлементов в технике, пер. $\begin{aligned} & \text { с англ., М., } 1982 . \text { В. Б. Андреев. } \\ & \text { ЭМДЕНА УРАВНЕНИЕ - нелинейное обыкновенное }\end{aligned}$ дифференциальное уравнение 2-го порядка



\[

\frac{d^{2} y}{d x^{2}}+\frac{2}{x} \frac{d y}{d x}+y^{\alpha}=0

\]



или, в самосопряженной форме,



\[

\frac{d}{d x}\left(x^{2} \frac{d y}{d x}\right)+x^{2} y^{\alpha}=0

\]



где $\alpha>0, \alpha \neq 1$, - константа. Точка $x=0$ является для Э. у. особой. Заменой переменной $x=1 / \xi$ уравнение (1) приводится к виду



\[

\frac{d^{2} y}{d \xi^{2}}+\frac{y^{\alpha}}{\xi^{4}}=0

\]



а заменой $y=\eta / x-$ к виду



\[

\frac{d^{2} \eta}{d x^{2}}+\frac{\eta^{\alpha}}{x^{\alpha-1}}=0

\]



После замены переменных



\[

x=e^{-t}, y=e^{\mu t} u, \mu=2 /(\alpha-1),

\]



и последующего понижения порядка подстановкой $u^{\prime}=v(u)$ получается уравнение 1-го порядка



\[

\frac{d v}{d u}=-(2 \mu-1)-\frac{\mu(\mu-1) u+\mu^{\alpha}}{v} .

\]



У равнение (1) было получено Р. Эмденом [1] в связи с изучением условий равновесия политропного газового шара; эта задача сводится к задаче существования $\mathrm{y}$ уравнения (1) с начальными условиями $y(0)=1$, $y^{\prime}(0)=0$ решения, определенного на нек-ром отрезке $\left[0, x_{0}\right], 0<x_{0}<\infty$, и обладающего свойствами:



\[

y(x)>0 \text { при } 0 \leqslant x<x_{0}, y\left(x_{0}\right)=0 .

\]



Иногда уравнение (1) наз. также уравнением Л е їн а - Э м д е н $\mathbf{\text { а. }}$



Более общими, чем Э. у., являются у р а в н е н и е $\Phi$ a y л e p a



\[

\frac{d}{d x}\left(x^{2} \frac{d y}{d x}\right)+x^{\lambda} y^{\alpha}=0, \lambda>0, \alpha>0,

\]



б уравневие Әмдена-Фа уле р а



\[

\frac{d}{d x}\left(x^{\rho} \frac{d y}{d x}\right) \pm x^{\lambda} y^{\alpha}=0

\]



где $\rho, \lambda, \alpha \neq 1$ - действительные параметры. Как частный случай это уравнение включает у р а в в е н и е Т о м а с а - Ф е р м и



\[

\frac{d^{2} y}{d x^{2}}=\frac{y^{3 / 2}}{\sqrt{x}}

\]



возникающее при изучении распределения әлектронов в атоме. Если $\rho \neq 1$, то уравнение (2) заменоӥ переменных может быть преобразоваво к виду



\[

\frac{d^{2} w}{d s^{2}} \pm s^{\sigma} w^{\alpha}=0

\]



Имеются различные результаты качественного и асимптотич. исследования решений уравнения Эмдена Фаулера (см., напр., [2], [3]). Подробно изучалось также ур а внение типа Әмдена-Ф а ул е р а



\[

\frac{d^{2} y}{d x^{2}}+a(x)|y|^{\alpha} \operatorname{sign} y=0

\]



(см. о нем и его аналоге n-го порядка в [4]). Jum.: [1] E m d en R., Gaskugeln, Lpz.- B., 1907; [2] С а н с о н е Д ж., Обыкновенные пифференциальные уравне-

ния, пер. с итал., т. 2, М. 1954; [3].Б л л м а Р Р. теория устойчивости решений дифференциальных уравнений, пер. с англ., М., 1954; [4] К и г у р а д з е И. Т., Некоторые сингулярные краевые задачи для обыкновенных диффференциальных



уравнений, тб́., 1975 .



ӘМПИРИЧЕСКОЕ РАСПРЕДЕЛЕНИЕ, $p$ а с п $p$ eд е ле и е вы б о р к и, - распределение вероятностей, к-рое определяется по выборке для оценивания истинного распределения. Пусть результаты наблюдений $X_{1}, \ldots, X_{n}$ - взаимно независимые и одинаково распределенные случайные величины с функцией распределения $\mathscr{F}(x)$ и пусть $X_{(1)}<X_{(2)}<\ldots<X_{(n)}-$ coотвөтствующий вариационный ряд. р а с п р е д е л ени е м, соответствующим $X_{1}, \ldots$, $X_{n}$, наз. дискретное распределение, приписывающеө каждому значению $X_{k}$ вероятность $1 / n$. Функция Э. $p$. $\mathscr{F}_{n}(x)$ наз. э мп и р и ч е с к о й ф у н ц И е й р а cп р е д л е н и я, является ступенчатой функцией со скачками, кратными $1 / n$, в точках, определяемых величинами $X_{(1)}, \ldots$, $X_{(n)}$ :



\[

\hat{\mathscr{J}}_{n}(x)=\left\{\begin{array}{l}

0, x \leqslant X_{(1)}, \\

\frac{k}{n}, X_{(k)}<x \leqslant X_{(k+1)}, \quad 1 \leqslant k \leqslant n-1, \\

1, x>X_{(n)} .

\end{array}\right.

\]



При фиикспрованных значениях $X_{1}, \ldots, X_{n}$ функция $\hat{\mathscr{F}}_{n}(x)$ обладает всеми свойствами обычной функции распределения. При каждом фиксированном действительном $x$ функция $\hat{\mathscr{T}}_{n}(x)$ является случайной величиной как функция ${ }^{\prime} X_{1}, \ldots, X_{n}$. Таким образом, Э. р., соответствующее выборке $X_{1}, \ldots, X_{n}$, задается семейством случайных величин $\widehat{\mathscr{F}}_{n}(x)$, зависящих от действительного параметра $x$. При әтом для фиксированного $x$



И



\begin{center}

\includegraphics[max width=\textwidth]{2023_07_09_1c0a5893169b96c905ddg-1}

\end{center}



\[

\mathrm{P}\left\{\widehat{\mathscr{F}}_{n}(x)=\frac{k}{n}\right\}=C_{n}^{k}[\mathscr{F}(x)]^{k}[1-\mathscr{F}(x)]^{n-k} .

\]



В соответствии с законом больших чисел $\hat{\mathscr{F}}_{n}(x) \stackrel{\mathrm{P}}{\longrightarrow}$ $\longrightarrow F(x), n \longrightarrow \infty$, при каждом $x$. Әто означает, что $\hat{\mathscr{F}}_{n}(x)$ - несмещенная и состоятельная оценка функции распределения $\mathscr{F}(x)$. Функция Э. p. равномерно по $x$ сходится с вероятностью $1 \mathrm{k} \mathscr{J}(x)$ при $n \longrightarrow \infty$ или, если



\[

\begin{aligned}

& D_{n}=\sup _{x}\left|\hat{\mathscr{F}}_{n}(x)-\mathscr{F}(x)\right|, \\

& \mathrm{P}\left\{\lim _{n \rightarrow \infty} D_{n}=0\right\}=1

\end{aligned}

\]



(теорема $\Gamma$ ливенко- Кантелли).



Величина $D_{n}$ служит мерой близости $\hat{\mathscr{T}}_{n}(x)$ к $\mathscr{F}(x)$. А. Н. Колмогоров (1933) нашел предельное распределение: для непрерывной функции $\mathscr{F}(x)$



$\lim _{n \rightarrow \infty}\left\{\sqrt{n} D_{n}<z\right\}=K(z)=\sum_{n=-\infty}^{+\infty}(-1)^{k} e^{-2 k^{2} z^{2}}, z>0$.



Если $\mathscr{F}(x)$ неизвестна, то для проверки гипотезы о том, что эта функция есть заданная непрерывная функция $\mathscr{F}_{0}(x)$, применяются критерии, основанные на статистиках типа $D_{n}$ (см. Колмогорова критериц̆, Колмогорова - Смирнова критерий, Непараметрические меmoды статистики).



Моменты и любые другие характеристики Э. р. наз. вы б о р п ны м и (э м п и р и ч е с ки и и) моментами и характеристиками, напр.:



\[

\bar{X}=\frac{1}{n} \sum_{k=1}^{n} X_{k}-\text { выборочное среднее, }

\]





\end{document}